\documentclass[11pt,a4paper]{article}

% ---------- Font ve dil ----------
\usepackage{fontspec}
\setmainfont{DejaVu Sans}
\newfontfamily\turkishfont{DejaVu Sans}

% ---------- Sayfa düzeni ----------
\usepackage[a4paper,margin=2.4cm]{geometry}
\usepackage{fancyhdr}
\usepackage{titlesec}
\usepackage{enumitem}
\usepackage{booktabs}
\usepackage{array}
\usepackage{longtable}
\usepackage{xcolor}
\usepackage{tcolorbox}
\tcbuselibrary{skins,breakable}
\usepackage{hyperref}
\usepackage{listings}
\usepackage{parskip}

% ---------- Renkler ----------
\definecolor{codebg}{RGB}{245,246,248}
\definecolor{codeframe}{RGB}{210,214,220}
\definecolor{jsblue}{RGB}{0,92,197}
\definecolor{jsgreen}{RGB}{60,140,60}
\definecolor{jsgray}{RGB}{110,110,110}
\definecolor{noteorange}{RGB}{176,96,0}
\definecolor{noteorangebg}{RGB}{255,244,229}
\definecolor{titleblue}{RGB}{18,52,86}

% ---------- Başlık biçimleri ----------
\titleformat{\section}
  {\Large\bfseries\color{titleblue}}{\thesection}{0.6em}{}
\titleformat{\subsection}
  {\large\bfseries\color{titleblue}}{\thesubsection}{0.6em}{}

% ---------- Kod bloğu stili ----------
\lstdefinelanguage{JavaScriptTR}{
  keywords={var, let, const, function, return, if, else, for, while, true, false, typeof, new},
  keywordstyle=\color{jsblue}\bfseries,
  ndkeywords={console, log, Math, Number, parseFloat},
  ndkeywordstyle=\color{jsgreen},
  comment=[l]{//},
  commentstyle=\color{jsgray}\itshape,
  string=[b]",
  string=[b]',
  stringstyle=\color{noteorange},
  morestring=[b]`,
  sensitive=true
}

\lstset{
  language=JavaScriptTR,
  basicstyle=\ttfamily\small,
  backgroundcolor=\color{codebg},
  frame=single,
  rulecolor=\color{codeframe},
  framerule=0.4pt,
  breaklines=true,
  breakatwhitespace=false,
  showstringspaces=false,
  tabsize=2,
  columns=fullflexible,
  xleftmargin=8pt,
  xrightmargin=8pt,
  aboveskip=8pt,
  belowskip=8pt
}

% ---------- Dikkat / Not kutusu ----------
\newtcolorbox{dikkat}{
  breakable, enhanced,
  colback=noteorangebg, colframe=noteorange,
  boxrule=0.6pt, arc=2pt, left=6pt, right=6pt, top=4pt, bottom=4pt,
  title={\textbf{Dikkat}},
  fonttitle=\bfseries\color{noteorange},
  coltitle=noteorange
}

% ---------- Üstbilgi/altbilgi ----------
\pagestyle{fancy}
\setlength{\headheight}{14pt}
\fancyhf{}
\renewcommand{\headrulewidth}{0.3pt}
\fancyhead[L]{JavaScript Başlangıç Notları}
\fancyhead[R]{\thepage}

\hypersetup{
  colorlinks=true,
  linkcolor=titleblue,
  urlcolor=titleblue,
  pdftitle={JavaScript Başlangıç Notları},
  pdfauthor={Ayşenur}
}

\renewcommand{\contentsname}{İçindekiler}

\begin{document}

% ================= KAPAK =================
\begin{titlepage}
\centering
\vspace*{3.5cm}
{\Huge\bfseries\color{titleblue} JavaScript Başlangıç Notları}\\[0.8cm]
{\Large Değişkenler, Stringler, Sayılar, Diziler ve Nesneler}\\[2.5cm]
{\large Ayşenur}\\[0.3cm]
{\large 3 Ağustos 2026}\\[3cm]
{\normalsize Bu belge, \texttt{JavaScriptBaslangic} GitHub reposundaki\\
\texttt{3\_js\_kodlama} klasöründe yazılan pratik kodların ders notu haline getirilmiş halidir.}
\vfill
\end{titlepage}

\tableofcontents
\newpage

% ================= GİRİŞ =================
\section*{Klasör Yapısı}
\addcontentsline{toc}{section}{Klasör Yapısı}

\begin{lstlisting}[language={},basicstyle=\ttfamily\small,frame=single,backgroundcolor=\color{codebg}]
3_js_kodlama/
|-- index.html     -> script6.js'i sayfaya baglayan basit HTML dosyasi
|-- script.js       -> Degiskenler, veri tipleri, string temelleri
|-- script2.js      -> String uygulamalari (URL/slug ornegi)
|-- script3.js      -> Sayilar (Number, Math) ve diziler
|-- script4.js      -> Dizi metotlari
|-- script5.js      -> Dizi uygulamasi + nesneler (objects)
`-- script6.js      -> Ic ice nesne/dizi uygulamasi (siparis sistemi)
\end{lstlisting}

Konsol çıktılarını görmek için \texttt{index.html} dosyasını tarayıcıda açıp geliştirici araçları (F12) $\rightarrow$ Console sekmesine bakman yeterli.

\newpage

% ================= SCRIPT.JS =================
\section{script.js --- Değişkenler, Veri Tipleri, String Temelleri}

\subsection*{Değişken Tanımlama}
\texttt{var} anahtar kelimesiyle değişken tanımlandı (\texttt{let} de aynı işi görür, aralarındaki fark ileride kapsam/scope konusunda önemli olacak).

\begin{lstlisting}
var a = 5000;
a = 6000; // var ile tanimlanan degiskenin degeri sonradan degistirilebilir
\end{lstlisting}

\textbf{Değişken isimlendirme kuralları:} İsimler rakamla başlayamaz, JS'in ayrılmış komut isimleri (\texttt{var}, \texttt{function} vb.) değişken adı olarak kullanılamaz.

\subsection*{Veri Tipleri ve \texttt{typeof}}

\begin{lstlisting}
var urunAdi = 'Iphone 16';   // string
var urunfyat2 = 7000;         // number
console.log(typeof urunAdi);  // "string"
\end{lstlisting}

\begin{dikkat}
İki string \texttt{+} ile birleştirildiğinde matematiksel toplama değil, \textbf{concatenation (birleştirme)} yapılır:
\begin{lstlisting}
var sayi1 = '10';
var sayi2 = '20';
console.log(sayi1 + sayi2); // "1020" olur, 30 degil!
\end{lstlisting}
\end{dikkat}

\textbf{\texttt{undefined}:} Değer atanmamış bir değişkenin hem değeri hem tipi \texttt{undefined} döner.

\subsection*{Uygulama --- Öğrenci Not Ortalaması}
İki öğrencinin doğum yılından yaşını hesaplama (\texttt{2026 - doğumYılı}) ve üç sınav notunun ortalamasını alıp 50 barajına göre başarı durumunu (\texttt{true}/\texttt{false}) belirleme. Bu, karşılaştırma operatörlerinin (\texttt{>=}) boolean sonuç ürettiğini gösteren güzel bir örnek.

\subsection*{String'e Index ile Erişim ve Template String}

\begin{lstlisting}
var ad = 'Aysenur';
console.log(ad[0]); // "A"

var mesaj = `Benim adim ${ad}`; // template string (backtick)
\end{lstlisting}

\subsection*{String Metotları}

\begin{center}
\begin{tabular}{@{}p{4cm}p{9.5cm}@{}}
\toprule
\textbf{Metot} & \textbf{Ne yapar} \\
\midrule
\texttt{toUpperCase()} & Tüm harfleri büyütür \\
\texttt{length} & Karakter sayısını verir \\
\texttt{slice(start, end)} & Belirtilen aralığı keser; negatif index sağdan sayar \\
\texttt{substring(start, end)} & slice'a benzer, negatif index desteklemez \\
\texttt{replace(eski, yeni)} & İlk eşleşmeyi değiştirir \\
\texttt{trim()} & Baştaki/sondaki boşlukları siler \\
\texttt{indexOf(metin)} & Metnin bulunduğu index'i verir, yoksa -1 \\
\texttt{split(ayraç)} & Stringi diziye böler \\
\bottomrule
\end{tabular}
\end{center}

\begin{dikkat}
Dosyada \texttt{kursAdi.split(" ")[-1]} ile son elemanı almaya çalışılmış. JavaScript'te Python'daki gibi \textbf{negatif index desteği yoktur}, bu yüzden \texttt{[-1]} çalışmaz ve \texttt{undefined} döner. Son elemanı almak için \texttt{split(" ").at(-1)} ya da \texttt{split(" ")[dizi.length - 1]} kullanılmalı.
\end{dikkat}

\newpage

% ================= SCRIPT2.JS =================
\section{script2.js --- String Uygulamaları}

Gerçek bir senaryo: bir URL ve kurs adından ``slug'' (URL'de kullanılabilir kısa isim) üretme.

\begin{lstlisting}
let url = "https://www.sadikturan.com";
let kursAdi = "Komple Web Gelistirme Kursu";

var kacKelime = kursAdi.split(" ");     // kelimelere ayirma
var urlBulma = url.split(":");           // protokolu ayirma ("https")
var kucuk = kursAdi.toLowerCase();       // kucuk harfe cevirme
var yeni = kucuk.replaceAll(" ", "-");   // bosluklari tire ile degistirme

var mesaj = `${url}/${yeni}`;
console.log(mesaj); // https://www.sadikturan.com/komple-web-gelistirme-kursu
\end{lstlisting}

Bu dosyada öğrenilen ana fikir: string metotları zincirleme kullanılarak gerçek dünyada işe yarayan bir dönüşüm (slug oluşturma) yapılabiliyor.

\newpage

% ================= SCRIPT3.JS =================
\section{script3.js --- Sayılar ve Diziler}

\subsection*{Tip Dönüştürme ve Biçimlendirme}

\begin{lstlisting}
sonuc = Number(sonuc);       // string -> number
sonuc = parseFloat("20.2");   // ondalikli stringi sayiya cevirir

let sayi = 10.58946;
sayi.toPrecision(6); // toplam 6 basamak gosterir
sayi.toFixed(5);      // virgulden sonra 5 basamak gosterir
\end{lstlisting}

\subsection*{Math Nesnesi}

\begin{center}
\begin{tabular}{@{}p{4cm}p{9.5cm}@{}}
\toprule
\textbf{Fonksiyon} & \textbf{Ne yapar} \\
\midrule
\texttt{Math.round(x)} & En yakın tam sayıya yuvarlar \\
\texttt{Math.ceil(x)} & Yukarı yuvarlar \\
\texttt{Math.floor(x)} & Aşağı yuvarlar \\
\texttt{Math.sqrt(x)} & Karekök alır \\
\texttt{Math.pow(x, y)} & Üs alma ($x^y$) \\
\texttt{Math.min(...)} & En küçük değeri bulur \\
\texttt{Math.random()} & 0 ile 1 arasında rastgele sayı üretir \\
\bottomrule
\end{tabular}
\end{center}

\begin{dikkat}
\texttt{Math.random() * 10} $\rightarrow$ 0-10 arası, \texttt{Math.random() * 100 + 1} $\rightarrow$ 1-101 arası sayı üretmek için kullanılan genel kalıp.
\end{dikkat}

\subsection*{Diziler (Array) ve İç İçe Diziler}

\begin{lstlisting}
let urunler  = ["Iphone 15", "Iphone 16", "Iphone 17"];
let fiyatlar = [50000, 60000, 70000];

let urun11 = ["Iphone 15", 50000, ["gold", "silver", "black"]];
console.log(urun11[2][2]); // "black" -> ic dizinin 3. elemani
\end{lstlisting}

\newpage

% ================= SCRIPT4.JS =================
\section{script4.js --- Dizi Metotları}

Bir öğrenci listesi üzerinden dizilerde ekleme/çıkarma/arama işlemleri denenmiş:

\begin{center}
\begin{tabular}{@{}p{5.2cm}p{8.3cm}@{}}
\toprule
\textbf{Metot} & \textbf{Ne yapar} \\
\midrule
\texttt{length} & Eleman sayısı \\
\texttt{toString()} & Diziyi virgülle ayrılmış stringe çevirir \\
\texttt{join(' ')} & Diziyi belirtilen ayraçla stringe çevirir \\
\texttt{pop()} & Son elemanı \textbf{siler ve döndürür} \\
\texttt{shift()} & İlk elemanı \textbf{siler ve döndürür} \\
\texttt{push(x)} & Sona eleman \textbf{ekler}, dizinin yeni uzunluğunu döndürür \\
\texttt{unshift(x)} & Başa eleman \textbf{ekler}, dizinin yeni uzunluğunu döndürür \\
\texttt{indexOf(x)} & Elemanın index'ini bulur \\
\texttt{includes(x)} & Eleman var mı diye \texttt{true/false} döner \\
\texttt{splice(b, s, ...y)} & Belirtilen konumdan eleman siler/ekler \\
\bottomrule
\end{tabular}
\end{center}

\begin{dikkat}
\texttt{push()} ve \texttt{unshift()} eklenen elemanı değil, \textbf{dizinin yeni uzunluğunu (length)} döndürür. Kodda \texttt{sonuc = ogrenciler.push("Sena")} satırından sonra \texttt{console.log(sonuc)} bir sayı (dizinin yeni uzunluğu) basar, "Sena" değil --- bu JS'e yeni başlayanların en çok şaşırdığı noktalardan biri.
\end{dikkat}

\newpage

% ================= SCRIPT5.JS =================
\section{script5.js --- Dizi Uygulaması + Nesneler (Objects)}

İlk kısımda meyve dizisi üzerinde \texttt{push}/\texttt{pop} pratiği yapılmış.

\subsection*{Nesneler (Object) --- Key/Value Yapısı}

\begin{lstlisting}
let kullanici1 = {
  "ad": "Sadik",
  "soyad": "Turan",
  "yas": 40,
  "adres": {
    "sehir": "kocaeli",
    "ilce": "izmit"
  },
  "hobileler ": ["sinema", "spor"]
};
\end{lstlisting}

\subsection*{Nesne Değerlerine Erişim --- İki Yöntem}

\begin{lstlisting}
kullanici1["ad"];        // koseli parantez (bracket) notasyonu
kullanici1.adres.ilce;   // nokta (dot) notasyonu, ic ice nesnelerde de calisir
\end{lstlisting}

\begin{dikkat}
\texttt{"hobileler "} anahtarının sonunda fazladan bir boşluk var (yazım hatası). Bu yüzden \texttt{kullanici1.hobileler} çalışmaz, sadece \texttt{kullanici1["hobileler "]} (boşlukla birlikte) çalışır. Nesne anahtarlarında boşluk/özel karakter olursa dot notation kullanılamaz, bracket notation zorunlu hale gelir --- bu satır bunun canlı bir örneği.

Ayrıca \texttt{kullanicilar = [kullanici1, kullanici2];} satırında değişken \texttt{var}/\texttt{let} olmadan tanımlanmış. Bu, JS'te ``implicit global'' (örtük global değişken) oluşturur ve genelde önerilmez --- ileride \texttt{"use strict"} modunda bu satır hataya bile sebep olabilir.
\end{dikkat}

\newpage

% ================= SCRIPT6.JS =================
\section{script6.js --- Nesne Uygulaması: Sipariş Sistemi}

Gerçekçi bir e-ticaret sipariş verisi üzerinden iç içe (nested) nesne ve dizi kullanımı pekiştirilmiş.

\begin{lstlisting}
let siparis1 = {
  siparis_id: 101,
  satin_alinanlar: [           // dizi icinde nesneler (birden fazla urun)
    { urun_id: 5, urun_basligi: "Iphone 16 Pro", urun_fiyat: 75000 },
    { urun_id: 6, urun_basligi: "Iphone 16 Pro Max", urun_fiyat: 85000 }
  ],
  musteri: { musteri_id: 12 }
};
\end{lstlisting}

\texttt{siparis2} içinde ise \texttt{satin\_alinanlar} bir dizi değil, tek bir nesne --- yani aynı alan farklı sipariş kayıtlarında farklı veri yapısına (dizi vs. tekil nesne) sahip olabiliyor, bu da erişim şeklini değiştiriyor:

\begin{lstlisting}
// siparis1: dizi oldugu icin index ile erisim + toplama
var siparis_1toplam = (siparis1.satin_alinanlar[0].urun_fiyat
                      + siparis1.satin_alinanlar[1].urun_fiyat) * 1.2;

// siparis2: tekil nesne oldugu icin dogrudan alan adiyla erisim
var siparis_2toplam = siparis2.satin_alinanlar.urun_fiyat * 1.2;
\end{lstlisting}

\texttt{* 1.2} çarpanı \%20 KDV eklemek için kullanılmış. Bu dosyanın öğrettiği en önemli fikir: gerçek dünya verisi (API'den gelen JSON gibi) genelde bu şekilde iç içe nesne ve dizilerin karışımından oluşur, bu yüzden veriye erişmeden önce yapısının (dizi mi, nesne mi) bilinmesi gerekir.

\newpage

% ================= ÖZET =================
\section{Genel Özet --- Öğrenilen Konular}

\begin{itemize}[leftmargin=1.2em]
\item Değişken tanımlama (\texttt{var}, \texttt{let}) ve isimlendirme kuralları
\item Veri tipleri: \texttt{string}, \texttt{number}, \texttt{boolean}, \texttt{undefined} ve \texttt{typeof} operatörü
\item String birleştirme (\texttt{+}) ile matematiksel toplama arasındaki fark
\item Template literal (backtick) kullanımı
\item String metotları: \texttt{toUpperCase}, \texttt{slice}, \texttt{substring}, \texttt{replace}, \texttt{trim}, \texttt{indexOf}, \texttt{split}
\item Sayı dönüştürme (\texttt{Number}, \texttt{parseFloat}) ve biçimlendirme (\texttt{toFixed}, \texttt{toPrecision})
\item \texttt{Math} nesnesi fonksiyonları (\texttt{round}, \texttt{ceil}, \texttt{floor}, \texttt{sqrt}, \texttt{pow}, \texttt{min}, \texttt{random})
\item Diziler: tanımlama, index ile erişim, iç içe diziler
\item Dizi metotları: \texttt{push}, \texttt{pop}, \texttt{shift}, \texttt{unshift}, \texttt{indexOf}, \texttt{includes}, \texttt{splice}, \texttt{join}
\item Nesneler (objects): key-value yapısı, iç içe nesneler, dot/bracket notation
\item Gerçekçi veri modelleme: dizi + nesne kombinasyonlarıyla (sipariş sistemi) çalışma
\end{itemize}

\section{Tekrar Ederken Dikkat Edilecek Noktalar}

\begin{itemize}[leftmargin=1.2em]
\item \texttt{split(" ")[-1]} JS'te çalışmaz --- negatif index yoktur, \texttt{.at(-1)} kullanılmalı.
\item \texttt{push()} / \texttt{unshift()} eklenen elemanı değil, dizinin \textbf{yeni uzunluğunu} döndürür.
\item Nesne anahtarında boşluk/özel karakter varsa \textbf{sadece bracket notation} (\texttt{obj["anahtar adı"]}) çalışır, dot notation çalışmaz.
\item Değişkeni \texttt{var}/\texttt{let} olmadan atamak örtük global değişken oluşturur, kaçınılmalı.
\item Aynı alan (\texttt{satin\_alinanlar} gibi) farklı kayıtlarda dizi ya da tekil nesne olabilir --- erişimden önce veri yapısı kontrol edilmeli.
\end{itemize}

\end{document}
